\documentclass[11pt]{article}

\usepackage[margin=1in]{geometry}
\usepackage{amsmath}
\usepackage{amssymb}
\usepackage{booktabs}
\usepackage{longtable}

\title{DIS Fit Model Mapping Reference}
\author{a1n\_analysis}
\date{\today}

\begin{document}

\maketitle

\section{Purpose}

This note documents the Deep Inelastic Scattering (DIS) fit models currently
wired into the \texttt{a1n\_analysis} pipeline through
\texttt{src/dis\_fit\_models.py}.  For each selectable model key, this note
records the functional form, parameter ordering, and the partial-derivative
mapping used for uncertainty propagation through \texttt{fit\_error(...)}.

For all models with an explicit $x^\alpha$ prefactor, the live code enforces
\[
\alpha \ge 10^{-6},
\]
so the fit cannot diverge as $x \to 0$.

\section{Selector Keys}

\begin{longtable}{llll}
\toprule
Model key & Python function & Display name & Beta index \\
\midrule
\texttt{fullx}        & \texttt{g1f1\_quad\_fullx\_DIS}   & Full-x Quadratic      & 4 \\
\texttt{quad\_alpha}  & \texttt{g1f1\_quad\_alpha\_DIS}   & Power-Law Quadratic   & 4 \\
\texttt{cubic\_alpha} & \texttt{g1f1\_cubic\_alpha\_DIS}  & Power-Law Cubic       & 5 \\
\texttt{quad2}        & \texttt{g1f1\_quad2\_DIS}         & Constrained Quadratic & 3 \\
\texttt{quad}         & \texttt{g1f1\_quad\_DIS}          & Quadratic             & 3 \\
\texttt{cubic}        & \texttt{g1f1\_cubic\_DIS}         & Cubic                 & 4 \\
\bottomrule
\end{longtable}

Backward-compatible aliases accepted by the code:
\begin{itemize}
\item \texttt{quad\_new} $\rightarrow$ \texttt{quad\_alpha}
\item \texttt{cubic\_new} $\rightarrow$ \texttt{cubic\_alpha}
\item \texttt{quadnew} $\rightarrow$ \texttt{quad\_alpha}
\item \texttt{cubicnew} $\rightarrow$ \texttt{cubic\_alpha}
\end{itemize}

The special selector \texttt{DIS\_FIT\_MODEL = "all"} does not define a new
functional form.  It runs all supported models and selects the one minimizing
\[
\left| \chi^2_{\mathrm{red}} - 1 \right|.
\]

\section{Model-by-Model Mapping}

\subsection{\texttt{fullx}}

Parameter order:
\[
(\alpha, a, b, c, \beta, d, x_0, \sigma).
\]

Functional form:
\begin{equation}
\frac{g_1}{F_1}(x,Q^2)
=
x^\alpha
\left[
a + bx + cx^2 + d \exp\!\left(-\frac{1}{2}\left(\frac{x-x_0}{\sigma}\right)^2\right)
\right]
\left(1 + \frac{\beta}{Q^2}\right).
\end{equation}

Partials used by the code:
\begin{align*}
\alpha &\rightarrow \texttt{partial\_alpha\_fullx}, &
a      &\rightarrow \texttt{partial\_a\_fullx}, \\
b      &\rightarrow \texttt{partial\_b\_fullx}, &
c      &\rightarrow \texttt{partial\_c\_fullx}, \\
\beta  &\rightarrow \texttt{partial\_beta\_fullx}, &
d      &\rightarrow \texttt{partial\_d\_fullx}, \\
x_0    &\rightarrow \texttt{partial\_x0\_fullx}, &
\sigma &\rightarrow \texttt{partial\_sigma\_fullx}.
\end{align*}

\subsection{\texttt{quad\_alpha}}

Parameter order:
\[
(\alpha, a, b, c, \beta).
\]

Functional form:
\begin{equation}
\frac{g_1}{F_1}(x,Q^2)
=
x^\alpha (a + bx + cx^2)\left(1 + \frac{\beta}{Q^2}\right).
\end{equation}

Partials used by the code:
\begin{align*}
\alpha &\rightarrow \texttt{partial\_alpha\_quad\_alpha}, &
a      &\rightarrow \texttt{partial\_a\_quad\_alpha}, \\
b      &\rightarrow \texttt{partial\_b\_quad\_alpha}, &
c      &\rightarrow \texttt{partial\_c\_quad\_alpha}, \\
\beta  &\rightarrow \texttt{partial\_beta\_quad\_alpha}.
\end{align*}

\subsection{\texttt{cubic\_alpha}}

Parameter order:
\[
(\alpha, a, b, c, d, \beta).
\]

Functional form:
\begin{equation}
\frac{g_1}{F_1}(x,Q^2)
=
x^\alpha (a + bx + cx^2 + dx^3)\left(1 + \frac{\beta}{Q^2}\right).
\end{equation}

Partials used by the code:
\begin{align*}
\alpha &\rightarrow \texttt{partial\_alpha\_cubic\_new}, &
a      &\rightarrow \texttt{partial\_a\_cubic\_new}, \\
b      &\rightarrow \texttt{partial\_b\_cubic\_new}, &
c      &\rightarrow \texttt{partial\_c\_cubic\_new}, \\
d      &\rightarrow \texttt{partial\_d\_cubic\_new}, &
\beta  &\rightarrow \texttt{partial\_beta\_cubic\_new}.
\end{align*}

\subsection{\texttt{quad2}}

Parameter order:
\[
(x_0, y_0, c, \beta).
\]

Functional form:
\begin{equation}
\frac{g_1}{F_1}(x,Q^2)
=
\left[c(x-x_0)^2 + y_0\right]\left(1+\frac{\beta}{Q^2}\right).
\end{equation}

Partials used by the code:
\begin{align*}
x_0   &\rightarrow \texttt{partial\_x0}, &
y_0   &\rightarrow \texttt{partial\_y0}, \\
c     &\rightarrow \texttt{partial\_c4}, &
\beta &\rightarrow \texttt{partial\_beta4}.
\end{align*}

\subsection{\texttt{quad}}

Parameter order:
\[
(a, b, c, \beta).
\]

Functional form:
\begin{equation}
\frac{g_1}{F_1}(x,Q^2)
=
(a + bx + cx^2)\left(1+\frac{\beta}{Q^2}\right).
\end{equation}

Partials used by the code:
\begin{align*}
a     &\rightarrow \texttt{partial\_a2}, &
b     &\rightarrow \texttt{partial\_b2}, \\
c     &\rightarrow \texttt{partial\_c2}, &
\beta &\rightarrow \texttt{partial\_beta2}.
\end{align*}

\subsection{\texttt{cubic}}

Parameter order:
\[
(a, b, c, d, \beta).
\]

Functional form:
\begin{equation}
\frac{g_1}{F_1}(x,Q^2)
=
(a + bx + cx^2 + dx^3)\left(1+\frac{\beta}{Q^2}\right).
\end{equation}

Partials used by the code:
\begin{align*}
a     &\rightarrow \texttt{partial\_a3}, &
b     &\rightarrow \texttt{partial\_b3}, \\
c     &\rightarrow \texttt{partial\_c3}, &
d     &\rightarrow \texttt{partial\_d3}, \\
\beta &\rightarrow \texttt{partial\_beta3}.
\end{align*}

\section{Legacy Partial Naming Note}

The comment in \texttt{src/functions.py}
\begin{quote}
\texttt{\# partial functions for quadratic and cubic forms} \\
\texttt{\# 2 means for quad form, 3 means for cubic form}
\end{quote}
describes only a legacy subset of derivative names:
\begin{itemize}
\item \texttt{*2} for the plain quadratic model,
\item \texttt{*3} for the plain cubic model,
\item \texttt{*4} for the constrained quadratic model,
\item \texttt{*\_fullx} for the full-x quadratic model,
\item \texttt{*\_new} for the original power-law quadratic notation.
\end{itemize}

The active model selection path does not infer derivative mappings from that
comment.  Instead, \texttt{src/dis\_fit\_models.py} explicitly binds each model
key to its fit function, parameter order, and partial list.

\end{document}
